\documentclass[11pt]{article}
\usepackage[margin=1in]{geometry}
\usepackage{graphicx}
\usepackage{booktabs}
\usepackage{amsmath,amssymb}
\usepackage{natbib}
\usepackage{hyperref}
\usepackage{xcolor}
\usepackage{enumitem}
\graphicspath{{../outputs/figures/}{../results/figures/}}
\newcommand{\todo}[1]{\textcolor{red}{[TODO: #1]}}
\newcommand{\coauthor}[1]{\textcolor{blue!60!black}{[COAUTHOR: #1]}}
\title{\textbf{Degrees of Separation --- Next-Stage Research Agenda}\\[4pt]
\large A forward-looking companion to the static paper:\\
the platform, the frontier, and the two leverage points}
\author{Sora Yng \and \coauthor{prospective advisors/collaborators --- see closing}}
\date{Agenda companion to \texttt{main.tex}. Synthesised from the analysis record; no new analysis.}

\begin{document}
\maketitle

\begin{abstract}
This agenda is the forward-looking companion to the static paper (\texttt{paper/main.tex}). It \textbf{supersedes}
the prior research-agenda material (the project proposals, \texttt{degrees\_of\_separation\_proposal\{,\_v2\}.md},
and the agenda fragments in earlier drafts), which predate two changes that strengthen the base: the
\emph{revealed-placement} reframe of the second axis, and the placement-\emph{identification} layer that
establishes earnings as a validated-but-incomplete placement proxy with two characterised failure modes. The
static paper is now a \emph{platform} --- a measured, structured, cross-validated, mechanistically-explained wedge
--- and the frontier is to relax the four things the platform's placement axis still is: a single scalar at the
institution grain, US-only, median-level, and disconnected from student beliefs. Three thrusts relax these in
order of access cost: \textbf{(A)} cross-national external validity, available now on public data; \textbf{(B)} the
multi-dimensional, within-field placement gap, the flagship, gated on Revelio resume data, into which the prior
residual-mechanism, tail, and dynamic questions all converge; and \textbf{(C)} a behavioural layer linking the
measured gap to student subjective expectations, whose descriptive and quasi-experimental tiers run now on public
data and whose flagship RCT is gated on a co-author. Two leverage points define the programme: a Revelio
subscription unlocks the within-field, multi-dimensional, and tail frontier and the placement half of the dynamic
question (B; the prestige-temporal half needs a separate dynamic-SpringRank upgrade, edge density its binding
constraint); a subjective-expectations collaboration unlocks the flagship behavioural and welfare contribution (C).
\end{abstract}

% =====================================================================
\section{Where the project stands: the platform}\label{sec:platform}
The static paper establishes four findings, and they should be read as a \emph{platform}, not a set of claims to
defend. \textbf{(1) Measurement:} the academic-reputation--placement gap,
$\mathrm{gap}_f = 1 - \mathrm{Spearman}_i(\text{prestige}_{i,f},\,\text{placement}_{i,f})$ within field across
institutions, is a measurable, reliable object (only 15\% of fields inside a perfect-agreement-plus-noise null,
65\% above; artifact channel $R^2=0.20$; a per-field reliability filter). \textbf{(2) Structure:} it is
discipline-structured --- a precision-weighted multilevel decomposition places $\sim$30\% of cross-field gap
variance between broad discipline areas (ICC $0.30$; measurement-corrected $0.45$), with an integrated end
(Computing, Mathematics/Statistics, Economics-adjacent social science, Engineering) and a decoupled end (Health,
Arts, the natural sciences). \textbf{(3) A validated-but-incomplete placement proxy:} earnings tracks an
independent earnings source (PSEO vs Scorecard, institution pay Spearman $+0.94$; field gap $+0.68$) and non-wage
placement (occupational prestige, vertical match) at $+0.65$ to $+0.84$ across horizons and sources.
\textbf{(4) One identified mechanism:} occupational licensing cleanly moves the gap ($b_{\text{licensure}}=+0.66$),
and the mechanism is identified --- in regulated, locally-employed fields pay is set by \emph{setting}, not school
(as licensure rises, the incremental wage-variance explained by school prestige beyond geography falls,
$\mathrm{corr}(\text{licensure},\text{incremental-prestige }R^2)=-0.59$; Nursing $\approx 0$).

\paragraph{The placement axis, precisely.} The reframe matters for the agenda because it scopes the open
questions. The second axis is not ``employer reputation''; it is \emph{revealed labour-market placement} ---
where graduates land --- and earnings is its validated-but-incomplete proxy, with two \emph{characterised} failure
modes rather than vague ``limitations'': a cross-field \emph{proxy-timing} effect (in high-deferral fields
early-career earnings understates terminal placement; the status-residual tracks the graduate-degree share at
Pearson $+0.59$, distinct from the null academic-absorption channel and not touching the within-field gap,
$\mathrm{corr}(\text{gap},\text{deferral})=+0.02$), and the \emph{licensing-setting} effect of finding (4). The
platform is therefore not ``a gap we measured'' but ``a gap we measured, structured, cross-validated, and partly
mechanised --- whose placement axis we have \emph{identified} as a scalar proxy with known where-it-fails.''

% =====================================================================
\section{Organising logic of the frontier}\label{sec:logic}
The platform's placement axis is, deliberately, four things at once: \textbf{(a)} a single scalar (earnings) at the
institution grain; \textbf{(b)} US-only; \textbf{(c)} median-level and descriptive; and \textbf{(d)} disconnected
from the student beliefs that rankings actually move. Each frontier thrust relaxes one or more of these. Make
placement \emph{multi-dimensional and within-field} (relaxes a, and turns c from a level into a structure);
make it \emph{cross-national} (relaxes b, available now on public data); make it \emph{causal/tail-resolved}
(relaxes c fully, gated on administrative data); and make it \emph{behavioural} (relaxes d; public descriptive and
quasi-experimental tiers run now, only the flagship RCT gated on a co-author).

\paragraph{The dynamic question folds into Thrust B, but its binding constraint is edge density.} The prior agenda
treated the prestige--earnings lead--lag as its own thrust; the field-level pilot gave only two noisy periods and
could not identify a direction. The reason is \emph{not}, as the earlier framing implied, that prestige is static.
A prestige \emph{time series} is constructible from the project's existing ORCID hiring-network rebuild --- rebuild
the network by appointment-year window (ORCID's dated PhD$\to$first-faculty-appointment edges) and run SpringRank
per window to obtain $\text{prestige}_f(t)$ --- so it is an \emph{extension} of a pipeline already in the project,
not missing data. The binding constraint is \emph{edge density}, and it is worse at the temporal grain: splitting
the already-sparse, lossy hiring-edge set (the source of the $\rho\approx0.77$ reconstruction ceiling) across time
windows leaves each window too thin for a stable per-window SpringRank --- exactly why the pilot had two noisy
periods. The placement side (earnings over time: Scorecard cohorts, or Revelio) is adequate; the bottleneck is on
the \emph{prestige-network} side. The dynamic question therefore folds into Thrust B --- it needs the
institution-level placement panel \emph{and} a time-varying prestige series --- and resolving it needs one of three
upgrades on the prestige side: \textbf{(a)} a denser hiring-edge source (ORCID $+$ faculty directories $+$ Revelio
academic transitions); \textbf{(b)} a temporally-pooled \emph{dynamic} SpringRank with cross-window shrinkage /
state-space smoothing to stabilise sparse-window estimates --- a slowly-evolving-graph-recovery-under-edge-sparsity
problem in the spirit of minimax VAR-graph recovery; or \textbf{(c)} coarser windows (resolution traded for
density). We do not list the lead--lag separately, but flag edge density, not data availability, as its real
obstacle.

% =====================================================================
\section{Thrust A --- Cross-national external validity (immediate, public)}\label{sec:thrustA}
\paragraph{What and why.} The single analytical ceiling-raiser available \emph{before} any co-author or new
proprietary data is to ask whether the integrated/decoupled typology is a US institution or a feature of how
prestige and placement relate in general. This is external validity, and it is buildable now on public data.

\paragraph{Estimand.} For country $k$, recompute the field gap $\mathrm{gap}^{k}_f$ and its CIP-2 (or local
equivalent) cluster means; the test statistic is the cross-country rank correlation of the field-level typology,
\[
\tau_{\,US,k} = \mathrm{Spearman}_c\!\left(\overline{\mathrm{gap}}^{\,US}_c,\ \overline{\mathrm{gap}}^{\,k}_c\right),
\]
over discipline clusters $c$. A high $\tau$ says the integrated--decoupled ordering is a general property of the
prestige$\leftrightarrow$placement map, not a US artifact; a low $\tau$ localises it and is itself informative.

\paragraph{Data.} \emph{UK --- LEO} (Longitudinal Education Outcomes): provider $\times$ subject graduate earnings,
administratively linked, and crucially published \emph{region-reweighted} --- a built-in geography control that
directly addresses the setting-driven mechanism of finding (4), against a UK research-prestige hierarchy (REF or a
UK faculty-hiring network where constructible). \emph{Europe --- EUROGRADUATE}: a comparative graduate survey
across $\sim$18 countries that elicits \emph{self-reported horizontal skills-match} --- a direct, individual-level
placement-\emph{match} micro-measure, i.e.\ a non-wage placement dimension observed natively rather than proxied by
wages. EUROGRADUATE thus also previews the multi-dimensional placement of Thrust B at the survey grain.

\paragraph{Feasibility.} Low risk, public, no IRB, no subscription. The deliverable is a cross-national replication
table of the typology plus the $\tau_{US,k}$ statistic, and it is the natural first manuscript-strengthening
extension before the gated thrusts.

% =====================================================================
\section{Thrust B --- The multi-dimensional, within-field placement gap (flagship; Revelio-gated)}\label{sec:thrustB}
This is the programme's flagship, and it is where the prior agenda's separate residual-mechanism, tail, and dynamic
thrusts \emph{converge}, because they require the \emph{same} object: institution $\times$ field placement resolved
to occupation, firm, and time. One data unlock --- Revelio --- delivers the placement side of all of them; the
dynamic item (v) additionally needs a time-varying prestige series (\S\ref{sec:logic}).

\paragraph{(i) The within-field non-wage placement gap.} The direct analogue of the paper's gap, on a non-wage
placement dimension:
\[
\mathrm{gap}^{\mathrm{occ}}_f = 1 - \mathrm{Spearman}_i\!\left(\text{prestige}_{i,f},\ \text{occ-prestige}_{i,f}\right),
\qquad
\text{occ-prestige}_{i,f} = \sum_o P(o\mid i,f)\,\pi_o,
\]
where $P(o\mid i,f)$ is the occupation distribution of field-$f$ graduates from institution $i$ and $\pi_o$ a
survey occupational-prestige score. The static paper can only build the field-level version (one occupation
distribution per field); the within-field, across-institution version --- the true test of whether the typology is
a wage phenomenon or a general placement phenomenon --- needs institution $\times$ field occupation flows. This is
the estimand already specified in the placement-identification layer.

\paragraph{(ii) Sorting versus within-occupation pricing.} A law-of-total-variance decomposition separates the two
channels by which prestige could raise placement:
\[
\mathrm{Var}\!\left(\text{wage}\mid i,f\right)
= \underbrace{\mathrm{Var}_o\!\left(\mathbb{E}[\text{wage}\mid o]\right)}_{\text{between-occupation sorting}}
+ \underbrace{\mathbb{E}_o\!\left[\mathrm{Var}(\text{wage}\mid o)\right]}_{\text{within-occupation pricing}} .
\]
The estimand is the share of the prestige--placement relationship that runs through \emph{which occupations}
graduates enter (sorting) versus \emph{pay within occupation} (pricing) --- testing the employer-learning reading
of finding (4) directly.

\paragraph{(iii) Employer-prestige / firm-quality-weighted placement.} Re-weight placement by the prestige or
selectivity of the \emph{hiring firm} --- a second hierarchy estimable as a SpringRank over the firm-to-firm
hiring network, or as AKM-style firm pay premia. This makes ``where graduates land'' an employer-identity object,
the labour-market analogue of the faculty-hiring network that defines AR.

\paragraph{(iv) Career trajectory and promotion.} With individual transitions over time, placement becomes a
\emph{slope}, not a level: seniority and promotion trajectories test whether decoupled-field prestige buys a
steeper path even when it does not buy a higher starting wage.

\paragraph{(v) Institution-level dynamic lead--lag (the folded-in dynamic question).} The lead--lag becomes
identifiable at the institution grain, but it needs \emph{both} halves: the Revelio panel supplies the time-varying
\emph{placement} side, while the time-varying \emph{prestige} side must come from an ORCID-windowed dynamic
SpringRank (\S\ref{sec:logic}), whose binding constraint is edge density, not data availability. Revelio unlocks the
placement half; the prestige-temporal half needs the dynamic-SpringRank / denser-edge upgrade --- the study the
field-level pilot could not support.

\paragraph{(vi) Tail-resolved brand premium.} Resume data observe top-firm and top-occupation placement
\emph{rates}, not just the median --- the elite tail where \citet{chetty2023diversifying} locate the brand's
causal payoff and which the paper's median earnings cannot see. This converts the paper's ruled-out
``generic-brand'' alternative from a median-data limitation into a measurable quantity.

\paragraph{Data and access.} Revelio Labs' resume-and-employment-record panel carries the required dictionary:
occupation, firm, seniority, skills, geography, and individual job-to-job transitions --- exactly the six
coordinates the platform's scalar proxy omits. Access path: the Revelio research programme
(\href{mailto:research@reveliolabs.com}{research@reveliolabs.com}; academic engagements associated with
D.~Dorn and F.~Schoner) \coauthor{confirm current contact and academic-access terms before approach}. \emph{This
is why Revelio access is the programme's single highest-value unlock:} one subscription turns the within-field,
multi-dimensional, and tail frontier from ``specified'' into ``estimable,'' and supplies the placement half of the
dynamic frontier --- whose prestige-temporal half needs the separate dynamic-SpringRank upgrade of \S\ref{sec:logic}.

% =====================================================================
\section{Thrust C --- The behavioural layer (co-author-gated; NHB scope)}\label{sec:thrustC}
The platform measures a structure in \emph{outcomes}; rankings act on \emph{beliefs}. Closing that loop is the
boldest contribution and the natural locus of a subjective-expectations collaboration.

\paragraph{Three tiers by data requirement.} The behavioural layer does \emph{not} have to begin with primary data
collection: two of its three tiers are public and runnable now, and only the flagship RCT needs fielding.

\emph{Tier 1 --- descriptive belief accuracy (public, no fielding; runnable now).} Public expectations panels ---
HSLS:09 (NCES; expected major and expected occupation at age 30, with realised outcomes) and NLSY97 --- let us test
whether students' expected major$\to$occupation/earnings systematically mis-estimates placement, and whether the
mis-estimation \emph{tracks the integrated/decoupled typology}:
$\mathrm{corr}_f\!\left(\overline{\text{belief error}}_f,\ \mathrm{gap}_f\right)$, joinable to the field-level gap.
Do students mis-price exactly the decoupled fields where prestige and placement diverge?

\emph{Tier 2 --- the information shock (public, quasi-experimental, no fielding).} The College Scorecard
field-of-study earnings rollout (institution-level 2015; field-of-study $\sim$2019) is a natural information
treatment in the Wiswall--Zafar spirit \citep{wiswall2015determinants}: with IPEDS CIP completions before/after,
does major choice respond to newly-public earnings information \emph{more in high-gap / decoupled fields}, treatment
intensity being the revealed-earnings surprise versus priors? Identification caveats are explicit (staggered rollout
timing; the field-of-study $\sim$2019 release confounded by COVID).

\emph{Tier 3 --- the flagship RCT (fielded; co-author $+$ IRB-gated).} The clean causal/welfare claim: elicit
subjective probabilistic beliefs over field $\times$ tier \emph{placement} --- not just earnings ---
\citep{manski2004measuring}, randomly provide the true gap/placement, and re-elicit or track intended choice
\citep{wiswall2015determinants}; with the generic-vs-field horse race re-run on \emph{beliefs} (does the belief
weight on the overall brand exceed that on field strength, the cognitive counterpart of $c_G=+0.45 \ge c_F=+0.40$
on outcomes?), this closes the loop to the welfare cost of field-agnostic rankings \citep{luca2013salience}.
Fielding can be a modest online-panel module (Prolific, or the USC Understanding America Study), not a captive
multi-year cohort --- but it is the one tier that strictly requires primary collection, needing IRB and a
subjective-expectations co-author. It is the tier that reaches \emph{Nature Human Behaviour} scope (causal $+$
welfare).

\paragraph{Gate.} Tiers 1--2 are public and need neither co-author nor IRB; only Tier 3 needs fielding, and even
that is a modest panel module, not a captive study. Thrust C thus \emph{starts} immediately on public data and
\emph{escalates} to the co-author-gated RCT --- the natural locus of a subjective-expectations collaboration.

% =====================================================================
\section{Sequencing, data gates, and the two leverage points}\label{sec:sequence}
\begin{table}[h]\centering
\caption{The three thrusts by timing, gate, and scope.}
\label{tab:thrusts}
\small
\begin{tabular}{p{2.9cm}p{2.1cm}p{2.6cm}p{2.4cm}p{3.4cm}}
\toprule
Thrust & Timing & Gate & Data & What it unlocks \\
\midrule
\textbf{A} Cross-national & now & none (public) & UK LEO; EUROGRADUATE & external validity of the typology
($\tau_{US,k}$); the pre-collaboration ceiling-raiser \\
\textbf{B} Within-field multi-dimensional (flagship) & on access & Revelio subscription & Revelio resume panel
(occupation, firm, time) & within-field / multi-dimensional / tail frontier $+$ the placement half of the dynamic question \\
\textbf{C} Behavioural & T1--2 now, T3 later & public (T1--2); co-author + IRB (T3 RCT) & HSLS:09 / NLSY97;
Scorecard rollout $+$ IPEDS; panel RCT & the belief--gap link and the welfare cost of rankings (NHB scope) \\
\bottomrule
\end{tabular}
\end{table}

The programme has \textbf{two leverage points}, and naming them is the point of this agenda. First, a
\textbf{Revelio subscription} unlocks the placement side of the whole of Thrust B: the within-field non-wage gap,
the sorting-vs-pricing decomposition, employer-prestige-weighted placement, trajectory, the placement half of the
institution-level dynamic lead--lag, and the tail-resolved brand premium are \emph{one} data unlock, not six
projects --- which is what makes it the highest-value single move (the dynamic item's prestige-temporal half needs
the dynamic-SpringRank / denser-edge upgrade of \S\ref{sec:logic}). Second, a \textbf{subjective-expectations collaboration} unlocks Thrust C and with it the behavioural
and welfare contribution that a purely outcome-side analysis cannot reach. Thrust A needs neither and should run
immediately as the external-validity strengthener. The static paper is the platform; A, B, and C are the
multi-year programme built on it, and the sequence is set by access cost, not by importance.

\bibliographystyle{plainnat}
\bibliography{refs}

\end{document}
